\CWHeader{Лабораторная работа \textnumero 2 \enquote{Поисковый робот}}

Необходимо написать поисковый робот --- компонент обкачки документов, используя
любой язык программирования.

\begin{itemize}
    \item Единственным аргументом поисковому роботу подаётся путь до
    yaml-конфига, содержащий:
    \begin{itemize}
        \item данные для базы данных в секции \texttt{db};
        \item данные для робота в секции \texttt{logic}: задержка между
        обкачкой страницы;
        \item любые другие данные, необходимые для реализации логики
        поискового робота.
    \end{itemize}
    \item Сохранять в базе данных (например, MongoDB) документы со следующими
    полями:
    \begin{itemize}
        \item url, нормализованный;
        \item \enquote{сырой} html-текст документа;
        \item название источника;
        \item дата обкачки документа в формате Unix time stamp.
    \end{itemize}
    \item Поисковый робот можно остановить в любой момент, и при повторном
    запуске робот должен начать с того документа, с которого он остановился.
    \item Периодически он должен уметь переобкачивать документы, которые уже
    есть в базе, но только в том случае, если они изменились.
\end{itemize}

\pagebreak
